\section{Measurement, Decoherence, and the Emergence of Classicality}
  \label{sec:classical-limit}

  The structural origin of Bell-type correlations identified above naturally raises
  the question of their apparent absence in classical macroscopic phenomena.
  Everyday physical systems exhibit effectively local and statistically independent
  behavior, despite being ultimately composed of quantum constituents.
  Any satisfactory account of Bell violations must therefore explain not only their
  existence, but also their suppression in the classical limit.

  Within the present framework, this transition is understood as a change in the
  effective properties of the projection from underlying configurations to observable
  descriptions.
  While the projection is generally non-injective at microscopic scales, environmental
  interactions, coarse-graining, and limited observational resolution progressively
  reduce the number of distinct underlying configurations compatible with a given
  observable outcome.
  As a result, the effective mapping becomes approximately injective.

  Decoherence plays a central role in this process.
  Through continuous interaction with uncontrolled environmental degrees of freedom,
  interference between distinct underlying configurations is dynamically suppressed.
  From the perspective of observable descriptions, this suppression eliminates the
  global consistency constraints responsible for non-factorizable joint probabilities.
  Observable outcomes can then be treated as effectively independent, restoring the
  factorization required for classical statistical reasoning.

  In this regime, joint probabilities admit an approximate decomposition of the form
  \begin{equation}
    P(a,b \mid A,B) \approx P(a \mid A)\,P(b \mid B),
  \end{equation}
  up to corrections that are exponentially small in the strength of environmental
  coupling or the degree of coarse-graining.
  Bell inequalities are therefore not violated in practice, not because the underlying
  structure has changed, but because the effective description no longer retains access
  to non-injective correlations.

  Importantly, this recovery of classical behavior does not require the introduction
  of additional collapse postulates or modifications of quantum dynamics.
  It follows from the same structural mechanism responsible for Bell violations,
  applied in a regime where the projection becomes effectively one-to-one.
  Classicality thus corresponds to a limit in which observable descriptions admit an
  approximately complete specification, rendering joint probabilities effectively
  factorizable.

  From this viewpoint, the quantum-to-classical transition reflects a change in
  descriptive accessibility rather than a change in fundamental physical laws.
  Bell-type correlations are suppressed when the information loss associated with
  non-injective projection is masked by decoherence and coarse-graining.
  The classical world emerges as a regime in which effective injectivity is restored,
  and with it, the validity of classical notions of separability and locality.
